\section{Rancangan Peraturan Menteri ESDM tentang Kuota PLTS Atap Berpasangan dengan Pensiun Dini PLTU}
\label{lamp:permen}

\begin{center}
\textbf{RANCANGAN}\\
\textbf{PERATURAN MENTERI ENERGI DAN SUMBER DAYA MINERAL}\\
\textbf{REPUBLIK INDONESIA}\\
\textbf{NOMOR \ldots{} TAHUN 2027}\\[4pt]
\textbf{TENTANG}\\[4pt]
\textbf{KUOTA PEMBANGKIT LISTRIK TENAGA SURYA ATAP BERPASANGAN DENGAN PERCEPATAN PENGAKHIRAN WAKTU OPERASI PEMBANGKIT LISTRIK TENAGA UAP}\\[6pt]
DENGAN RAHMAT TUHAN YANG MAHA ESA\\[4pt]
MENTERI ENERGI DAN SUMBER DAYA MINERAL REPUBLIK INDONESIA,
\end{center}

\noindent\textbf{Menimbang:}
\begin{enumerate}[label=\alph*.]
  \item bahwa percepatan pengakhiran waktu operasi PLTU memerlukan sumber energi pengganti dan kepastian pendanaan agar tidak membebani keuangan negara dan keandalan sistem.
  \item bahwa PLTS atap merupakan sumber energi pengganti yang dapat dipasang cepat dengan investasi yang tersebar di sektor publik dan swasta.
  \item bahwa kuota PLTS atap dan percepatan pengakhiran PLTU perlu ditetapkan secara berpasangan agar pengurangan penjualan listrik diimbangi penghematan biaya pembangkitan.
  \item bahwa berdasarkan pertimbangan pada huruf a, huruf b, dan huruf c, perlu menetapkan Peraturan Menteri tentang Kuota PLTS Atap Berpasangan dengan Percepatan Pengakhiran Waktu Operasi PLTU.
\end{enumerate}

\noindent\textbf{Mengingat:}
\begin{enumerate}
  \item Undang-Undang Nomor 30 Tahun 2009 tentang Ketenagalistrikan sebagaimana telah diubah.
  \item Undang-Undang Nomor 16 Tahun 2016 tentang Pengesahan Paris Agreement.
  \item Peraturan Presiden Nomor 112 Tahun 2022 tentang Percepatan Pengembangan Energi Terbarukan untuk Penyediaan Tenaga Listrik.
  \item Peraturan Menteri ESDM Nomor 2 Tahun 2024 tentang PLTS Atap yang Terhubung pada Jaringan Tenaga Listrik.
  \item Peraturan Menteri ESDM Nomor 10 Tahun 2025 tentang Peta Jalan Transisi Energi Sektor Ketenagalistrikan.
\end{enumerate}

\begin{center}\textbf{MEMUTUSKAN:}\end{center}
\noindent Menetapkan: PERATURAN MENTERI ENERGI DAN SUMBER DAYA MINERAL TENTANG KUOTA PLTS ATAP BERPASANGAN DENGAN PERCEPATAN PENGAKHIRAN WAKTU OPERASI PLTU.

\begin{center}\textbf{BAB I KETENTUAN UMUM}\end{center}
\noindent\textbf{Pasal 1.} Dalam Peraturan Menteri ini yang dimaksud dengan:
\begin{enumerate}
  \item Kuota Berpasangan adalah kuota PLTS atap pada suatu sistem yang ditetapkan dengan memperhitungkan kapasitas PLTU yang dijadwalkan mengakhiri operasi lebih awal pada sistem tersebut.
  \item PLTU Kandidat adalah PLTU yang masuk daftar urut percepatan pengakhiran operasi.
  \item Alat Hitung Bersama adalah perangkat perhitungan terbuka yang dipakai untuk menilai dampak usulan terhadap PT PLN (Persero), APBN, pelanggan, keandalan, emisi, dan kesehatan.
\end{enumerate}

\begin{center}\textbf{BAB II KUOTA BERPASANGAN}\end{center}
\noindent\textbf{Pasal 2.} (1) Kuota PLTS atap tahunan pada suatu sistem paling sedikit sama dengan kapasitas PLTU Kandidat yang dijadwalkan berhenti beroperasi pada sistem tersebut dalam tiga tahun berikutnya. (2) Kuota sebagaimana dimaksud pada ayat (1) tidak boleh lebih kecil dari kuota yang telah ditetapkan sebelumnya.

\noindent\textbf{Pasal 3.} (1) Instansi pemerintah dan BUMN wajib memasang PLTS atap pada gedung yang dimiliki atau dikelolanya sesuai target tahunan yang ditetapkan Menteri. (2) Untuk sistem Jawa-Madura-Bali, target sebagaimana dimaksud pada ayat (1) paling sedikit 200 MW per tahun mulai 2028.

\begin{center}\textbf{BAB III DAFTAR URUT DAN SYARAT PENUTUPAN}\end{center}
\noindent\textbf{Pasal 4.} (1) Daftar urut PLTU Kandidat disusun berdasarkan kriteria sebagaimana diatur dalam Peraturan Menteri ESDM Nomor 10 Tahun 2025 dengan tambahan pertimbangan jenis kontrak dan beban kesehatan per lokasi. (2) Penilaian dilakukan dengan Alat Hitung Bersama dan hasil agregatnya diumumkan.

\noindent\textbf{Pasal 5.} PLTU Kandidat dapat diakhiri operasinya apabila: (a) kompensasi telah disepakati dan paling sedikit 80\% pendanaannya telah terikat, (b) pembangkit pengganti telah beroperasi, (c) tambahan PLTS atap kumulatif pada sistem paling sedikit sama dengan kapasitas yang diakhiri, dan (d) program pelatihan ulang pekerja telah berjalan paling sedikit satu tahun.

\begin{center}\textbf{BAB IV PEMANTAUAN}\end{center}
\noindent\textbf{Pasal 6.} (1) Satuan Tugas Transisi Energi meninjau capaian kuota, penutupan, dan indikator transisi pada akhir setiap fase peta jalan. (2) PT PLN (Persero) menyampaikan data realisasi pembangkitan dan parameter kontrak untuk keperluan Alat Hitung Bersama setiap tahun.

\begin{center}\textbf{BAB V KETENTUAN PENUTUP}\end{center}
\noindent\textbf{Pasal 7.} Peraturan Menteri ini mulai berlaku pada tanggal diundangkan.

\section{Lembar Penilaian Kandidat Pensiun Dini}
\label{lamp:lembar}

Lembar ini diisi Satgas untuk setiap PLTU kandidat pada tahap 1 SOP. Angka diambil dari tab Kandidat PLTU, Indikator roadmap, dan Kriteria Permen 10/2025 di PowerShift. Contoh berikut diisi untuk Cirebon-1 dengan tahun pensiun 2030.

\begin{table}[H]
\centering
\small
\begin{tabularx}{\textwidth}{L{5.2cm} Y L{2.6cm}}
\toprule
Isian & Contoh Cirebon-1 & Sumber di PowerShift \\
\midrule
Pemilik dan jenis kontrak & IPP, PPA hingga Juli 2042 & Popup peta \\
Kapasitas dan usia pada tahun pensiun & 660 MW, 18 tahun & Popup peta \\
Tahun pensiun usulan & 2030 & Panel kebijakan \\
Kompensasi (estimasi ETM) & Rp7,6 triliun selama 2030-2040 & Tab Neraca \\
Neraca negara bila dipensiunkan sendiri & +Rp4,3 triliun (+Rp6,5 triliun per GW) & Tab Kandidat PLTU \\
Pembangkit pengganti yang dibutuhkan & 660 MW & Kartu angka \\
Beban kesehatan saat beroperasi & sekitar 910 kematian dini per tahun & Tab Kandidat PLTU \\
Pekerja langsung & sekitar 300 orang & Tab Indikator roadmap \\
Syarat 1: neraca negara per GW $\geq$ 0 & Terpenuhi & Tab Kandidat PLTU \\
Syarat 2: PLTS atap tambahan $\geq$ daya PLTU & Terpenuhi (sekitar 680 MWp pada 2030) & Tab Indikator roadmap \\
Syarat 3: dana luar terikat $\geq$ 80\% & Diisi saat negosiasi & Panel kebijakan \\
Syarat 4: pengganti beroperasi & Diisi saat penutupan & Tab Keandalan \\
Keputusan & Masuk daftar urut Fase 2 & \\
\bottomrule
\end{tabularx}
\end{table}

\noindent Langkah pengisian:
\begin{enumerate}
  \item Buka PowerShift, pilih skenario roadmap, lalu buka tab Kandidat PLTU dan pilih tahun pensiun usulan.
  \item Catat posisi PLTU pada daftar urut, neraca negara, dan beban kesehatannya.
  \item Tambahkan PLTU ke skenario, atur tahun pensiun dan porsi pendanaan di panel kebijakan.
  \item Baca status syarat pada tab Indikator roadmap dan tab Kriteria Permen 10/2025.
  \item Simpan skenario dengan tombol Bagikan dan lampirkan tautannya pada lembar penilaian.
\end{enumerate}

\section{Pemutakhiran Data dari Laporan Kelompok}
\label{lamp:pemutakhiran}

Kajian ini mengutamakan data yang sumbernya dapat diperiksa ulang dan cukup rinci untuk dimodelkan. Tabel berikut mencatat angka laporan kelompok yang dimutakhirkan.

\begin{table}[H]
\centering
\small
\begin{tabularx}{\textwidth}{L{3.6cm} L{3.8cm} Y}
\toprule
Laporan kelompok & Dokumen kajian & Alasan \\
\midrule
Target PLTS atap 2 GW pada 2028 & Kuota 1,85 GW Jawa-Madura-Bali hingga 2028 & Rincian wilayah dari sumber yang sama \cite{B-CNBC-ROOFTOP-2028}, sesuai cakupan kajian \\
\addlinespace
PLTS atap terpasang 772,9 MWp (November 2025) & 538 MWp (Juli 2025) sebagai titik awal model & Angka resmi KESDM yang dapat diperiksa \cite{B-KONTAN-ROOFTOP-2025}, laju tambahan dihitung dari periode yang sama \\
\addlinespace
Harga PLTS atap Rp15-20 juta per kWp (survei 2022) & Rp11-15 juta per kWp per segmen & Harga lebih baru \cite{B-DETIK-HARGA-PLTS} dan dibedakan antara rumah tangga, bisnis, industri, dan gedung pemerintah \\
\addlinespace
Valuasi pensiun dini per PLTU dari ETP & Estimasi ETM yang dikalibrasi ke refinancing Cirebon-1, ETP sebagai uji sensitivitas & Nilai ETP bergantung pada tahun pensiun versi ETP. Refinancing Cirebon-1 US\$250-300 juta adalah transaksi yang terverifikasi \cite{B-REUTERS-CIREBON-2022} \\
\addlinespace
Intensitas emisi rata-rata 790 g CO$_2$/kWh & Faktor emisi per unit PLTU & Dihitung per PLTU dari data GEM \cite{D-GEM-GCPT} \\
\addlinespace
Sekitar 30.000 pekerja sektor batu bara nasional & Pekerja langsung per PLTU yang dipensiunkan & Indikator dibuat per PLTU agar dapat dijadwalkan terhadap tanggal penutupan \cite{B-ECOBUSINESS-CIREBON} \\
\addlinespace
Target pensiun dini 1,7 GW dan alokasi US\$1,3 miliar & Tetap dipakai & Dikonfirmasi tanggapan CREA atas CIPP \cite{P-CREA-CIPP-2023}. Cakupan waktu dana belum pasti sehingga dipakai sebagai pembanding kasar \\
\bottomrule
\end{tabularx}
\end{table}

\section{Parameter Model}
\label{lamp:asumsi}

Status terverifikasi berarti angka dicocokkan ke dokumen asli. Status asumsi berarti pilihan tim kajian yang terdokumentasi. Kode sumber merujuk ke daftar pustaka dan halaman Sumber pada PowerShift.

{\small
\begin{longtable}{p{0.33\textwidth} p{0.2\textwidth} p{0.26\textwidth} p{0.11\textwidth}}
\toprule
Parameter & Nilai & Sumber & Status \\
\midrule
\endhead
Tahun dasar & 2025 tahun & \texttt{-\allowbreak{}} & asumsi \\
Beban puncak tertinggi Jawa-Madura-Bali & 28.077 MW & \texttt{B-\allowbreak{}MATAHARI-\allowbreak{}2024} & terverifikasi \\
Reserve margin Jawa-Madura-Bali 2025 & 0,375 fraksi & \texttt{B-\allowbreak{}PETROMINDO-\allowbreak{}RESERVE} & terverifikasi \\
Pertumbuhan beban dan penjualan & 0,0445 per tahun & \texttt{B-\allowbreak{}PETROMINDO-\allowbreak{}RESERVE} & terverifikasi \\
Tambahan daya mampu 2026-2028 & 1.000 MW/tahun & \texttt{B-\allowbreak{}PETROMINDO-\allowbreak{}RESERVE} & asumsi \\
Batas bawah reserve margin yang direkomendasikan & 0,33 fraksi & \texttt{B-\allowbreak{}PETROMINDO-\allowbreak{}RESERVE} & terverifikasi \\
Penjualan listrik PLN 2025 & 317,69 TWh & \texttt{B-\allowbreak{}CNBC-\allowbreak{}SALES-\allowbreak{}2025} & terverifikasi \\
Porsi penjualan Jawa-Bali & 0,695 fraksi & \texttt{D-\allowbreak{}PLN-\allowbreak{}STAT-\allowbreak{}2025} & terverifikasi \\
Biaya pokok penyediaan 2026 & 1.874 Rp/kWh & \texttt{B-\allowbreak{}KONTAN-\allowbreak{}SUBSIDI-\allowbreak{}2026} & terverifikasi \\
Kurs & 16.500 Rp/US\$ & \texttt{B-\allowbreak{}KONTAN-\allowbreak{}SUBSIDI-\allowbreak{}2026} & terverifikasi \\
Tingkat diskonto NPV & 0,08 per tahun & \texttt{assumption} & asumsi \\
Porsi kerugian PLN yang ditutup APBN & 1 fraksi & \texttt{assumption} & asumsi \\
Mulai tahun PLN menambah pembangkit untuk menjaga target cadangan & 2029 tahun & \texttt{B-\allowbreak{}LISTRIKID-\allowbreak{}RUPTL-\allowbreak{}JAMALI} & asumsi \\
Harga batu bara DMO untuk PLN & 70 US\$/ton & \texttt{B-\allowbreak{}KONTAN-\allowbreak{}DMO-\allowbreak{}2026} & terverifikasi \\
Kalori acuan harga DMO & 6.322 kcal/kg & \texttt{B-\allowbreak{}KONTAN-\allowbreak{}DMO-\allowbreak{}2026} & perlu cek \\
Biaya variabel non bahan bakar PLTU & 50 Rp/kWh & \texttt{assumption} & asumsi \\
Pembayaran kapasitas IPP (komponen A) & 175 US\$/kW/tahun & \texttt{P-\allowbreak{}IEEFA-\allowbreak{}2017} & asumsi \\
Biaya tetap O\&M PLTU milik PLN & 40 US\$/kW/tahun & \texttt{assumption} & asumsi \\
Nilai buku PLTU milik PLN tahun 2022 & 762.000 US\$/MW & \texttt{B-\allowbreak{}PELABUHANRATU-\allowbreak{}VALUASI} & perlu cek \\
Kompensasi pensiun dini, estimasi skema ETM & 1,05 miliar Rp/MW/tahun dipercepat & \texttt{B-\allowbreak{}REUTERS-\allowbreak{}CIREBON-\allowbreak{}2022} & terverifikasi \\
Kompensasi pensiun dini, klaim PLN & 18,18 miliar Rp/MW/tahun dipercepat & \texttt{B-\allowbreak{}PETROMINDO-\allowbreak{}CIREBON-\allowbreak{}BATAL} & terverifikasi \\
Pekerja langsung per MW & 0,4545 orang/MW & \texttt{B-\allowbreak{}ECOBUSINESS-\allowbreak{}CIREBON} & terverifikasi \\
Biaya kapasitas pembangkit pengganti & 120 US\$/kW/tahun & \texttt{assumption} & asumsi \\
Kematian dini rata-rata armada PLTU & 258,6 kematian/GW/tahun & \texttt{P-\allowbreak{}CREA-\allowbreak{}IESR-\allowbreak{}2023} & terverifikasi \\
Biaya kesehatan per kematian & 704.762 US\$/kematian & \texttt{P-\allowbreak{}CREA-\allowbreak{}IESR-\allowbreak{}2023} & terverifikasi \\
PLTS atap terpasang nasional Juli 2025 & 538 MWp & \texttt{B-\allowbreak{}KONTAN-\allowbreak{}ROOFTOP-\allowbreak{}2025} & terverifikasi \\
Porsi PLTS atap di Jawa-Bali & 0,8 fraksi & \texttt{assumption} & asumsi \\
Tambahan PLTS atap Jawa-Bali dengan aturan sekarang & 240 MWp/tahun & \texttt{B-\allowbreak{}KONTAN-\allowbreak{}ROOFTOP-\allowbreak{}2025} & asumsi \\
Kuota PLTS atap Jawa-Bali per tahun & 800 MW/tahun & \texttt{B-\allowbreak{}PVMAG-\allowbreak{}QUOTA} & asumsi \\
Koefisien kurva pangsa pasar vs balik modal & 0,3 1/tahun & \texttt{P-\allowbreak{}NREL-\allowbreak{}SOLARDS} & perlu cek \\
Umur PLTS atap & 25 tahun & \texttt{assumption} & asumsi \\
Lapangan kerja PLTS atap & 24,72 job-year/MW & \texttt{P-\allowbreak{}CEEW-\allowbreak{}2017} & terverifikasi \\
Keluaran puncak AC per kWp & 0,75 fraksi & \texttt{assumption} & asumsi \\
PLTS atap: Rumah tangga non subsidi & Rp15 jt/kWp, pakai sendiri 40\%, tarif Rp1.444,70 & \texttt{B-\allowbreak{}KOMPASTV-\allowbreak{}TARIF-\allowbreak{}2026} & asumsi \\
PLTS atap: Bisnis & Rp13 jt/kWp, pakai sendiri 80\%, tarif Rp1.224,83 & \texttt{D-\allowbreak{}PLN-\allowbreak{}STAT-\allowbreak{}2025} & asumsi \\
PLTS atap: Industri & Rp11 jt/kWp, pakai sendiri 90\%, tarif Rp1.085,39 & \texttt{D-\allowbreak{}PLN-\allowbreak{}STAT-\allowbreak{}2025} & asumsi \\
PLTS atap: Gedung pemerintah & Rp13 jt/kWp, pakai sendiri 80\%, tarif Rp1.583,12 & \texttt{D-\allowbreak{}PLN-\allowbreak{}STAT-\allowbreak{}2025} & asumsi \\
PLTS atap: Program rumah tangga subsidi & Rp12 jt/kWp, pakai sendiri 50\%, tarif Rp605,00 & \texttt{B-\allowbreak{}KOMPASTV-\allowbreak{}TARIF-\allowbreak{}2026} & asumsi \\
Lama respons pasar PLTS atap terhadap perubahan aturan & 4 tahun & \texttt{assumption} & asumsi \\
Subsidi listrik 2026 & 101.170 miliar Rp/tahun & \texttt{B-\allowbreak{}KONTAN-\allowbreak{}SUBSIDI-\allowbreak{}2026} & terverifikasi \\
Pelanggan subsidi 2026 & 44.880.000 pelanggan & \texttt{B-\allowbreak{}KONTAN-\allowbreak{}SUBSIDI-\allowbreak{}2026} & terverifikasi \\
Porsi subsidi di Jawa-Bali & 0,61 fraksi & \texttt{D-\allowbreak{}PLN-\allowbreak{}STAT-\allowbreak{}2025} & asumsi \\
Kuota PLTS atap Jawa-Madura-Bali hingga 2028 & 1.850 MWp & \texttt{B-\allowbreak{}CNBC-\allowbreak{}ROOFTOP-\allowbreak{}2028} & terverifikasi \\
Target pensiun dini nasional hingga 2040 (CIPP JETP) & 1,7 GW & \texttt{P-\allowbreak{}JETP-\allowbreak{}CIPP-\allowbreak{}2023} & terverifikasi \\
Alokasi pensiun dini CIPP JETP & 1,3 miliar US\$ & \texttt{P-\allowbreak{}JETP-\allowbreak{}CIPP-\allowbreak{}2023} & perlu cek \\
Pekerja terdampak tercakup pelatihan ulang atau perlindungan sosial & 1 fraksi & \texttt{L-\allowbreak{}LAPORAN-\allowbreak{}K17} & asumsi \\
\bottomrule
\end{longtable}

}

\section{Hitung Tangan}
\label{lamp:hitung}

Hitung tangan independen (\texttt{validation/cirebon1\_manual.py}) memakai skenario Cirebon-1 pensiun 2035 dengan dorongan PLTS atap dan kompensasi ETM. Seluruh komponen cocok dengan model dengan selisih di bawah 1\%.

\begin{table}[H]
\centering
\small
\begin{tabular}{lrrr}
\toprule
Komponen (Rp T, kecuali disebut) & 2030 & 2035 & 2040 \\
\midrule
Kapasitas PLTS atap (MWp) & 5.937 & 12.525 & 19.112 \\
Pendapatan PLN hilang & -4,80 & -10,54 & -16,29 \\
Kompensasi ekspor PLTS & -2,87 & -6,44 & -10,01 \\
Biaya energi dihemat & 3,03 & 6,75 & 10,55 \\
Pembayaran kapasitas dihemat & 0,00 & 1,91 & 1,91 \\
Porsi PLN kompensasi pensiun & 0,00 & -0,17 & -0,17 \\
Pembangkit pengganti & 0,00 & -1,31 & -1,31 \\
Neraca PLN & -4,64 & -9,80 & -15,32 \\
\bottomrule
\end{tabular}

\end{table}

\section{Daftar Unit PLTU}
\label{lamp:unit}

{\small
\begin{longtable}{p{0.22\textwidth} p{0.15\textwidth} r r l r r}
\toprule
PLTU & Unit & MW & Mulai & Tipe & Pensiun baseline & tCO$_2$/MWh \\
\midrule
\endhead
Adipala & Unit 1 & 660 & 2015 & PLN & 2055 & 0,94 \\
Batang (Central Java Power) & Unit 1 & 950 & 2022 & IPP & 2052 & 0,84 \\
Batang (Central Java Power) & Unit 2 & 950 & 2022 & IPP & 2052 & 0,84 \\
Celukan Bawang & Unit 1 & 142 & 2015 & IPP & 2045 & 1,15 \\
Celukan Bawang & Unit 2 & 142 & 2015 & IPP & 2045 & 1,15 \\
Celukan Bawang & Unit 3 & 142 & 2015 & IPP & 2045 & 1,15 \\
Cilacap Sumber & Phase I Unit 1 & 300 & 2006 & IPP & 2036 & 1,15 \\
Cilacap Sumber & Phase I Unit 2 & 300 & 2006 & IPP & 2036 & 1,15 \\
Cilacap Sumber & Phase II Unit 3 & 660 & 2016 & IPP & 2046 & 0,94 \\
Cilacap Sumber & Phase III Unit 4 & 1.000 & 2019 & IPP & 2049 & 0,84 \\
Cirebon & Unit 1 & 660 & 2012 & IPP & 2042 & 0,94 \\
Cirebon & Unit 2 & 924 & 2023 & IPP & 2053 & 0,84 \\
Indramayu & Unit 1 & 330 & 2010 & PLN & 2050 & 1,06 \\
Indramayu & Unit 2 & 330 & 2011 & PLN & 2051 & 1,06 \\
Indramayu & Unit 3 & 330 & 2011 & PLN & 2051 & 1,06 \\
Jawa-7 & Unit 1 & 991 & 2019 & IPP & 2049 & 0,84 \\
Jawa-7 & Unit 2 & 991 & 2020 & IPP & 2050 & 0,84 \\
Labuan & Unit 1 & 300 & 2009 & PLN & 2049 & 1,06 \\
Labuan & Unit 2 & 300 & 2010 & PLN & 2050 & 1,06 \\
Lontar & Unit 1 & 315 & 2011 & PLN & 2051 & 1,06 \\
Lontar & Unit 2 & 315 & 2012 & PLN & 2052 & 1,06 \\
Lontar & Unit 3 & 315 & 2012 & PLN & 2052 & 1,06 \\
Lontar & Unit 4 & 315 & 2024 & PLN & 2064 & 0,92 \\
Pacitan & Unit 1 & 315 & 2011 & PLN & 2051 & 1,06 \\
Pacitan & Unit 2 & 315 & 2011 & PLN & 2051 & 1,06 \\
Paiton & Unit 1 & 400 & 1994 & PLN & 2030 & 1,19 \\
Paiton & Unit 2 & 400 & 1993 & PLN & 2033 & 1,19 \\
Paiton Baru & Unit 09 & 660 & 2012 & PLN & 2052 & 0,97 \\
Paiton-1 & Unit 7 & 615 & 1999 & IPP & 2029 & 1,06 \\
Paiton-1 & Unit 8 & 615 & 1999 & IPP & 2029 & 1,06 \\
Paiton-2 & Unit 1 & 660 & 2000 & IPP & 2030 & 1,06 \\
Paiton-2 & Unit 2 & 660 & 2000 & IPP & 2030 & 1,06 \\
Paiton-3 & Unit 1 & 815 & 2012 & IPP & 2042 & 0,94 \\
Pelabuhan Ratu & Unit 1 & 350 & 2013 & PLN & 2045 & 1,06 \\
Pelabuhan Ratu & Unit 2 & 350 & 2013 & PLN & 2045 & 1,06 \\
Pelabuhan Ratu & Unit 3 & 350 & 2013 & PLN & 2045 & 1,06 \\
Rembang & Unit 1 & 315 & 2011 & PLN & 2051 & 1,06 \\
Rembang & Unit 2 & 315 & 2011 & PLN & 2051 & 1,06 \\
Serang & Unit 1 & 660 & 2017 & IPP & 2047 & 0,85 \\
Suralaya & Phase I Unit 1 & 400 & 1984 & PLN & 2030 & 1,19 \\
Suralaya & Phase I Unit 2 & 400 & 1985 & PLN & 2030 & 1,19 \\
Suralaya & Phase I Unit 3 & 400 & 1988 & PLN & 2030 & 1,19 \\
Suralaya & Phase I Unit 4 & 400 & 1989 & PLN & 2030 & 1,19 \\
Suralaya & Phase II Unit 1 & 600 & 1996 & PLN & 2036 & 1,10 \\
Suralaya & Phase II Unit 2 & 600 & 1997 & PLN & 2037 & 1,06 \\
Suralaya & Phase II Unit 3 & 600 & 1997 & PLN & 2037 & 1,06 \\
Suralaya & Unit 08 & 625 & 2011 & PLN & 2051 & 0,97 \\
Suralaya & Unit 09 & 1.000 & 2025 & IPP & 2055 & 0,84 \\
Suralaya & Unit 10 & 1.000 & 2025 & IPP & 2055 & 0,84 \\
Tanjung Awar-Awar & Unit 1 & 350 & 2012 & PLN & 2052 & 1,06 \\
Tanjung Awar-Awar & Unit 2 & 350 & 2013 & PLN & 2053 & 1,06 \\
Tanjung Jati B & Unit 1 & 660 & 2006 & IPP & 2036 & 1,04 \\
Tanjung Jati B & Unit 2 & 660 & 2006 & IPP & 2036 & 1,04 \\
Tanjung Jati B & Unit 3 & 660 & 2011 & IPP & 2041 & 0,96 \\
Tanjung Jati B & Unit 4 & 660 & 2012 & IPP & 2042 & 0,96 \\
Tanjung Jati B & Unit 5 & 1.000 & 2022 & IPP & 2047 & 0,84 \\
Tanjung Jati B & Unit 6 & 1.000 & 2022 & IPP & 2047 & 0,84 \\
\bottomrule
\end{longtable}

}
